\documentclass[11pt,letterpaper]{article}
\PassOptionsToPackage{hyphens}{url}
\usepackage[margin=.83in,top=.8in,bottom=.8in,headheight=15pt]{geometry}
\usepackage[T1]{fontenc}
\usepackage{lmodern,amsmath,amssymb,mathtools,bm,booktabs,array,graphicx,tikz}
\usetikzlibrary{arrows.meta,positioning,calc,decorations.pathreplacing,patterns}
\definecolor{ink}{HTML}{17344D}\definecolor{accent}{HTML}{087E8B}\definecolor{pale}{HTML}{EDF5F5}\definecolor{gold}{HTML}{AD6B19}
\providecommand{\dd}{\,\mathrm{d}}\providecommand{\R}{\mathbb{R}}
\providecommand{\E}{\mathbb{E}}\providecommand{\Var}{\operatorname{Var}}
\providecommand{\Cov}{\operatorname{Cov}}\providecommand{\diag}{\operatorname{diag}}
\providecommand{\norm}[1]{\left\lVert#1\right\rVert}
\providecommand{\abs}[1]{\left|#1\right|}
\usepackage{microtype,hyperref,fancyhdr,enumitem,amsthm,sectsty,needspace}
\allsectionsfont{\raggedright}
\hypersetup{colorlinks=true,linkcolor=ink,urlcolor=accent,pdfauthor={Zan Ahmad},pdftitle={Heat-flow operators, generators and finite spectral families / Review}}
\setlength{\parindent}{0pt}\setlength{\parskip}{6pt}\linespread{1.07}
\setlist{itemsep=3pt,topsep=4pt,leftmargin=*}
\emergencystretch=2em\clubpenalty=10000\widowpenalty=10000
\pagestyle{fancy}\fancyhf{}\fancyhead[L]{\small Computational Cardiology}\fancyhead[R]{\small Lecture 17}
\fancyfoot[L]{\footnotesize Zan Ahmad}\fancyfoot[R]{\footnotesize\thepage}\renewcommand{\headrulewidth}{.3pt}
\newtheorem{proposition}{Proposition}\newtheorem{theorem}{Theorem}\newtheorem{lemma}{Lemma}
\theoremstyle{definition}\newtheorem{definition}{Definition}\newtheorem{example}{Example}
\numberwithin{equation}{section}
\begin{document}
\section*{Heat-flow operators, generators and finite spectral families --- review sheet}
Use this after attempting the worked examples and practice. Explain each claim in words before calculating.

\section*{Key ideas}
\begin{itemize}
\item A function and its sampled vector are different objects.
\item Generators, flows and observables answer different questions.
\item A dataset split must match the desired generalization unit.
\item An exact synthetic reference allows targeted numerical checks.
\end{itemize}
\section*{Common errors and repairs}
\textbf{Calling any fixed-size regression a universally transferable operator.} State the tested family and discretizations.

\textbf{Splitting snapshots or nodes from one case across sets.} Group by the scientific unit intended to be unseen.

\section*{Retrieve without your notes}
\begin{itemize}
\item What does the zero-frequency heat mode do?
\item Does semigroup consistency establish biological validity?
\end{itemize}
\section*{Readiness check}
\begin{itemize}
\item Distinguish a field, its samples, a scalar observable, a generator and a solution map.
\item Specify an operator-learning problem with its domain, parameters, boundary conditions and prediction horizon.
\item Construct complete-case train/test regimes and identify unsupported distribution changes.
\item Check a simple diffusion operator against analytic and discrete references.
\end{itemize}
\end{document}
